\Subhead{4.2.1}{Резонаторы Фабри--Перо, конфокальный, концентрический и неустойчивый. Продольные и поперечные типы колебаний. Спектр, расходимость и добротность}

Оптическая мода воспроизводит пространственный профиль после полного обхода с комплексным множителем. В плоском резонаторе Фабри–Перо фазовое условие $2nk_0L+\phi_{r1}+\phi_{r2}=2\pi q$ даёт приблизительный межмодовый интервал $\Delta\nu_{\rm FSR}=c/(2n_gL)$ при однородном заполнении; групповая длина включает дисперсию среды и зеркал.

Для двух сферических зеркал $g_j=1-L/R_j$ и устойчивого параксиального семейства частоты можно записать
\begin{equation}
 \nu_{qmn}=\frac{c}{2nL}\left[q+\frac{m+n+1}{\pi}\arccos[\operatorname{sgn}(g_1)\sqrt{g_1g_2}]\right]
\end{equation}
с выбором ветви фазы обхода, соответствующей геометрии; знак/целую часть можно перераспределять в $q$. Дополнительный член — фаза Гуи. Вблизи вырожденных и граничных геометрий нужно следить за этой ветвью; выражение без её выбора не описывает все перестройки спектра.

Плоский резонатор $R_j=\infty$ находится на границе геометрической устойчивости и чувствителен к дифракции. Симметричный конфокальный $R_1=R_2=L$ имеет совпадающие фокусы и вырожденные поперечные семейства. Концентрический $R_1=R_2=L/2$ — граница; слегка смещённая внутрь области геометрия может быть устойчивой. В неустойчивом резонаторе лучи расширяются при обходах, но конечная апертура и усиление могут поддерживать генерирующую распределённую моду с выводом энергии через край: «неустойчивый» здесь не означает невозможность лазера.

Пусть энергия пассивной моды убывает $W=W_0\e^{-t/\tau_p}$. Тогда
\begin{equation}
 Q=\omega\tau_p=\frac{\nu}{\delta\nu},\qquad
 \delta\nu=\frac1{2\pi\tau_p}.
\end{equation}
$\delta\nu$ — полная ширина пассивного резонанса мощности, а амплитудный декремент $1/(2\tau_p)$. Если за время обхода $T_R$ энергия умножается на $\mathcal R<1$, $\tau_p=-T_R/\ln\mathcal R$; при малых потерях $\tau_p\simeq T_R/(1-\mathcal R)$. Финесс $\mathcal F=\Delta\nu_{\rm FSR}/\delta\nu$. Поперечные моды Эрмита–Гаусса различаются узлами и большей расходимостью; конечные зеркала дают зависящие от моды дифракционные потери.
